\documentclass[10pt]{article}
\usepackage[a4paper,landscape,top=12mm,bottom=12mm,left=13mm,right=13mm]{geometry}
\usepackage[T1]{fontenc}
\usepackage[utf8]{inputenc}
\usepackage[portuguese]{babel}
\usepackage{graphicx,xcolor,amsmath,microtype,booktabs}
\usepackage[hidelinks]{hyperref}
\pagestyle{empty}
\setlength{\parindent}{0pt}
\setlength{\parskip}{4pt}

\definecolor{tinta}{HTML}{272727}
\definecolor{cinza}{HTML}{4D4D4D}
\definecolor{claro}{HTML}{767676}
\definecolor{vermelho}{HTML}{B64342}
\definecolor{azul}{HTML}{0F4D92}
\definecolor{teal}{HTML}{42949E}
\color{tinta}

% cabeçalho de página de análise
\newcommand{\cab}[3]{%
  {\color{claro}\footnotesize\textsf{#1}}\\[-2pt]
  {\LARGE\textbf{#2}}\\[2pt]
  {\color{cinza}\small #3}\par\vspace{4pt}
  {\color{claro}\rule{\linewidth}{0.4pt}}\par\vspace{5pt}}

% rótulo de bloco no painel de texto
\newcommand{\bloco}[1]{\vspace{2pt}{\color{vermelho}\footnotesize\textbf{#1}}\par\vspace{-2pt}}
\newcommand{\ressalva}[1]{{\color{claro}\scriptsize #1}\par\vspace{1pt}}
\newcommand{\lit}[1]{\vspace{2pt}{\color{azul}\footnotesize\textbf{Na literatura}}\par\vspace{-2pt}{\color{cinza}\scriptsize #1}\par\vspace{1pt}}

% página: figura à esquerda, leitura à direita
\newcommand{\analise}[6]{%
  \cab{#1}{#2}{#3}
  \begin{minipage}[t]{0.588\linewidth}\vspace{0pt}
    \includegraphics[width=\linewidth]{#4}\par\vspace{6pt}
    {\color{claro}\rule{\linewidth}{0.4pt}}\par\vspace{4pt}
    {\scriptsize\setlength{\parskip}{0pt}\color{cinza}#6\par}
  \end{minipage}\hfill
  \begin{minipage}[t]{0.383\linewidth}\vspace{0pt}%
    \footnotesize\setlength{\parskip}{3pt}#5
  \end{minipage}
  \clearpage}

\newcommand{\analisec}[6]{%
  \cab{#1}{#2}{#3}
  \begin{minipage}[t]{0.588\linewidth}\vspace{0pt}
    \includegraphics[width=\linewidth]{#4}\par\vspace{6pt}
    {\color{claro}\rule{\linewidth}{0.4pt}}\par\vspace{4pt}
    {\scriptsize\setlength{\parskip}{0pt}\color{cinza}#6\par}
  \end{minipage}\hfill
  \begin{minipage}[t]{0.383\linewidth}\vspace{0pt}%
    \scriptsize\setlength{\parskip}{2.5pt}#5
  \end{minipage}
  \clearpage}


\newcommand{\analisea}[6]{%
  \cab{#1}{#2}{#3}
  \begin{minipage}[t]{0.515\linewidth}\vspace{0pt}
    \includegraphics[width=\linewidth]{#4}\par\vspace{6pt}
    {\color{claro}\rule{\linewidth}{0.4pt}}\par\vspace{4pt}
    {\scriptsize\setlength{\parskip}{0pt}\color{cinza}#6\par}
  \end{minipage}\hfill
  \begin{minipage}[t]{0.455\linewidth}\vspace{0pt}%
    \scriptsize\setlength{\parskip}{2.5pt}#5
  \end{minipage}
  \clearpage}



\newcommand{\legI}{\textbf{Fig. 1 | O corte de geração fotovoltaica no SIN é sistêmico, de razão energética e incide sobre toda a janela solar.} \textbf{a}, Perfil diurno médio da potência fotovoltaica das usinas com apuração de \textit{constrained-off} no Sistema Interligado Nacional. A área azul é a geração verificada e a área vermelha é a energia cortada, definida como \texttt{max(referência $-$ verificada, 0)} nos intervalos com restrição apurada; a soma define a potência que estaria disponível. Médias sobre todos os intervalos de 30 min do período; horário de Brasília, não corrigido para hora solar local. \textbf{b}, Taxa mensal de corte por subsistema, definida como energia cortada dividida pela energia disponível (verificada + cortada). \textbf{c}, Energia total cortada por razão de restrição, decomposta em origem sistêmica e local; rótulos indicam o total por razão. \textbf{d}, Distribuição espacial das usinas e conjuntos fotovoltaicos com apuração de \textit{constrained-off}; a cor codifica a taxa de corte acumulada e o tamanho, a capacidade instalada. A escala de cor satura no percentil 97 para preservar contraste. Contornos estaduais: IBGE. Dados: ONS, dados abertos, restrição por \textit{constrained-off} de usinas fotovoltaicas e fator de capacidade. Dados-fonte em \texttt{data/processed/fig1\_source\_*.csv}.}
\newcommand{\legII}{\textbf{Fig. 2 | A MMGD brasileira é majoritariamente sub-pixel no Sentinel-2, e a série de conexões carrega uma descontinuidade regulatória.} \textbf{a}, Distribuição acumulada da área do arranjo das instalações fotovoltaicas de micro e minigeração distribuída registradas na ANEEL. Linhas tracejadas marcam a área de um pixel de cada sensor; o rótulo indica a fração de instalações menor que um pixel. \textbf{b}, Fração de instalações que atinge o piso operacional de 4 pixels, por classe de consumo e sensor. \textbf{c}, Conexões mensais; a linha tracejada marca o fim da janela de transição da Lei 14.300/2022 e o marcador, o mês de pico. \textbf{d}, Participação da classe rural no total de instalações de cada unidade da federação. Contornos: IBGE. Dados: ANEEL, relação de empreendimentos de geração distribuída e informações técnicas fotovoltaicas. Dados-fonte em \texttt{data/processed/fig2\_source\_*.csv}.}
\newcommand{\legIII}{\textbf{Fig. 3 | Dois ativos solares em Minas Gerais diferem 3,0$\times$ na taxa de corte, e a diferença não se reproduz como padrão no parque.} \textbf{a}, Perfil diurno médio de potência, em pu da capacidade instalada, na janela comum aos dois ativos. Azul, geração verificada; vermelho, energia cortada. \textbf{b}, Custo Marginal de Operação do subsistema Sudeste ponderado pela energia, durante os intervalos de corte e durante a geração. \textbf{c}, Taxa de corte mensal dos dois ativos. \textbf{d}, Pixels Sentinel-2 ocupados pela área de módulos, comparados à instalação mediana de MMGD; escala logarítmica, área de módulos estimada a 5,12 m\textsuperscript{2}/kWp. \textbf{e}, Taxa de corte contra capacidade renovável variável num raio de 100 km, para os conjuntos com apuração completa na janela de referência da frota; os dois alvos em destaque. Dados: ONS (constrained-off, fator de capacidade, CMO, relacionamento conjunto-usina) e ANEEL (SIGA). Dados-fonte em \texttt{data/processed/fig3\_source\_*.csv}.}
\newcommand{\legIV}{\textbf{Fig. 4 | A queda do fator de capacidade de Sol do Cerrado acompanha a energização do entorno; o controle não mostra o mesmo.} \textbf{a}, Jaíba (MG). Área cinza, capacidade fotovoltaica acumulada em operação num raio de 100 km (eixo esquerdo); linha, fator de capacidade mensal verificado de Sol do Cerrado (eixo direito); faixa vermelha, parcela cortada apurada pelo ONS, disponível apenas a partir de abril de 2024 (linha tracejada). \textbf{b}, São Gonçalo do Abaeté (MG), com a mesma construção para Três Marias 3. \textbf{c}, Fator de capacidade mensal de Sol do Cerrado em 2023 e 2025, mês contra mês, para separar restrição de sazonalidade. \textbf{d}, Potência acumulada de micro e minigeração distribuída nos dois municípios; os 20,4 MW de Jaíba equivalem a 3,0\% da capacidade de Sol do Cerrado. Dados: ONS (fator de capacidade, constrained-off) e ANEEL (SIGA, MMGD). Dados-fonte em \texttt{data/processed/fig4\_source\_*.csv}.}
\newcommand{\legV}{\textbf{Fig. 5 | Normalizada pela irradiância medida sobre a usina, a perda de Sol do Cerrado permanece; o recurso solar não a explica.} \textbf{a}, Razão de desempenho mensal, definida como energia gerada por kW instalado dividida pela irradiância global horizontal diária; tracejado horizontal marca a média de cada ano. \textbf{b}, Irradiância global horizontal mensal sobre cada usina (linha cheia) e a mesma sob céu claro (tracejado). \textbf{c}, Decomposição da queda do fator de capacidade de Sol do Cerrado entre 2023 e 2025: observado em 2023, esperado em 2025 caso apenas a irradiância tivesse mudado, e observado em 2025. \textbf{d}, Fator de capacidade mensal contra irradiância mensal em Sol do Cerrado, por ano. Dados climáticos: NASA POWER (MERRA-2 + SYN1DEG), série diária no ponto de cada usina. Dados de energia: ONS. Dados-fonte em \texttt{data/processed/fig5\_source\_*.csv}.}
\newcommand{\legVI}{\textbf{Fig. 6 | A geração de referência do ONS é bem calibrada; o viés aparente da Fig. 5 vinha do contrafactual, não do operador.} \textbf{a}, Distribuição do viés mediano por usina, definido como a razão entre o potencial implicado pelo ONS (verificado + cortado) e o potencial previsto pela irradiância local usando a razão de desempenho calibrada nos dias sem restrição da própria usina; tracejado em 1 marca ausência de viés. \textbf{b}, Viés contra a intensidade mediana de corte nos dias restritos. \textbf{c}, Validação da calibração: fator de capacidade contra irradiância nos dias sem restrição de Sol do Cerrado, com a reta ajustada pela razão de desempenho limpa. \textbf{d}, Razão de desempenho anual de Sol do Cerrado contra o desempenho limpo calibrado; o rótulo interno indica o corte implícito de cada ano. Dados: ONS (fator de capacidade, constrained-off) e NASA POWER. Dados-fonte em \texttt{data/processed/fig6\_source\_*.csv}.}
\newcommand{\legVII}{\textbf{Fig. 7 | A construção de usinas fotovoltaicas é detectável e datável por Sentinel-2; a operação não é.} \textbf{a}, Razão entre o brilho no vermelho dentro do arranjo e num anel de controle de 1,5 a 3 km, por usina (linhas cinza) e mediana entre usinas (linha vermelha, faixa entre quartis), alinhadas pela data de entrada em operação do SIGA. \textbf{b}, Sentinel-2 em cor natural sobre o complexo Sol do Cerrado, em escala de refletância fixa entre as datas; círculos marcam as 17 usinas do SIGA. \textbf{c}, Comparação pareada da razão de brilho entre fases, uma linha por usina, com a mediana em destaque; \texttt{p} de Wilcoxon pareado. \textbf{d}, Diferença entre o ano de pico do brilho relativo e o ano de entrada em operação; barras em vermelho indicam erro de até um ano. Imagens: Sentinel-2 L2A (ESA), via catálogo STAC público. Registro: ANEEL/SIGA. Dados-fonte em \texttt{data/processed/fig7\_source\_*.csv}.}
\newcommand{\legVIII}{\textbf{Fig. 8 | O corte fotovoltaico por excesso de oferta coincide com térmica inflexível, mas apenas em carga baixa.} \textbf{a}, Perfil diurno médio nacional. Área azul, geração fotovoltaica verificada; área vermelha empilhada, energia cortada por razão energética (\texttt{ENE}); linhas, geração térmica separada em parcela econômica (por ordem de mérito, quando o CVU da usina é menor que o CMO) e parcela fora do mérito (geração total menos ordem de mérito). A anotação registra a carga líquida --- carga do SIN menos geração fotovoltaica --- na madrugada e no pico solar. \textbf{b}, Composição da geração térmica fora do mérito nas horas com e sem corte por razão energética, decomposta pelos motivos não-econômicos disjuntos do dicionário do ONS; rótulo indica o total. \textbf{c}, Correlação de Spearman entre corte por razão energética e geração térmica fora do mérito, estratificada por quintil de carga do SIN, restrita às horas de 9 h a 15 h; barras cinza indicam associação não significativa; abaixo do eixo, o corte médio de cada quintil. \textbf{d}, Controle negativo: taxa diária de corte por razão energética contra a energia armazenada do SIN, com a mediana por quartil de armazenamento. Dados: ONS, dados abertos --- restrição por \textit{constrained-off} fotovoltaico, geração térmica por motivo de despacho, curva de carga e energia armazenada por subsistema. Dados-fonte em \texttt{data/processed/fig8\_source\_*.csv}.}
\newcommand{\legIX}{\textbf{Fig. 9 | O piso térmico inflexível está onde o corte não está, e a coincidência nacional da Fig. 8 não sustenta leitura causal local.} \textbf{a}, Potência média nas horas solares (9 h--15 h) por subsistema: corte fotovoltaico por razão energética, geração térmica fora da ordem de mérito e a parcela desta que é inflexibilidade pura; abaixo do eixo, a razão entre o corte e o piso inflexível local. \textbf{b}, Correlação de Spearman entre corte por razão energética e térmica fora do mérito \textbf{do mesmo subsistema}, por quintil de carga do próprio subsistema; tracejado no limiar de magnitude declarado. \textbf{c}, A mesma correlação medida contra a térmica do próprio subsistema e contra a do outro subsistema com corte material. \textbf{d}, Distribuição acumulada do fluxo Nordeste$\rightarrow$Sudeste nas horas de corte alto (quartil superior do corte do Nordeste) e nas demais; tracejado no percentil 95 do fluxo, usado como referência empírica de teto por não haver limite operativo publicado. Dados: ONS, dados abertos --- \textit{constrained-off} fotovoltaico, geração térmica por motivo de despacho, curva de carga e intercâmbio nacional. Dados-fonte em \texttt{data/processed/fig9\_source\_*.csv}.}
\newcommand{\legX}{\textbf{Fig. 10 | A distância de rede até a demanda explica parte da alocação do corte que covariáveis locais não explicavam.} \textbf{a}, Grafo de transmissão em operação --- subestações e linhas da EPE, separadas em 500 kV ou mais e abaixo disso --- com as usinas fotovoltaicas da amostra posicionadas no seu ponto de conexão; a cor codifica a taxa de corte na janela de referência da frota e o tamanho, a capacidade instalada. Estrelas marcam as nove regiões metropolitanas usadas como centros de carga; a moldura cobre a região que contém a amostra, e 353 das 1.896 linhas do grafo ficam fora dela e não são desenhadas --- o recorte é de desenho, e todas as arestas participam do cálculo dos caminhos mínimos. \textbf{b}, Correlação de Spearman entre a taxa de corte e cada uma das oito covariáveis de rede pré-declaradas; em vermelho as que sobrevivem à correção de Holm para oito hipóteses. \textbf{c}, Taxa de corte contra a distância de rede ao centro de carga mais próximo, por subsistema, com a mediana por quartil de distância. \textbf{d}, As duas covariáveis sobreviventes medidas \textbf{dentro} de cada subsistema, para separar efeito de rede de efeito de subsistema. Dados: EPE, camadas de linhas de transmissão e subestações em operação; ONS, restrição por \textit{constrained-off} fotovoltaico e fator de capacidade. Dados-fonte em \texttt{data/processed/fig10\_source\_*.csv}.}
\newcommand{\legXI}{\textbf{Fig. 11 | O achado da Fig. 10 não se replica em eólica, e a dissociação espelha a origem do corte de cada fonte.} \textbf{a}, Correlação de Spearman entre a taxa de corte e as oito covariáveis de rede pré-declaradas na Fig. 10, medidas separadamente em fotovoltaica e eólica; a cor cheia marca as que sobrevivem à correção de Holm aplicada dentro de cada fonte, e o cinza as que não sobrevivem. \textbf{b}, Energia cortada decomposta por razão e origem da restrição, conforme classificação operativa do ONS, para cada fonte. \textbf{c}, As duas covariáveis-chave medidas apenas no Nordeste, onde as duas amostras coexistem, para separar tecnologia de geografia. \textbf{d}, Taxa de corte eólica contra o número de conjuntos no mesmo ponto de conexão, com a mediana por valor. Dados: ONS, restrição por \textit{constrained-off} eólico e fotovoltaico e fator de capacidade; EPE, linhas de transmissão e subestações em operação. Dados-fonte em \texttt{data/processed/fig11\_source\_*.csv}.}
\newcommand{\legXII}{\textbf{Fig. 12 | Dos quatro achados de rede da série, um passa nos três testes de robustez e o principal da Fig. 10 passa em apenas um.} \textbf{a}, Distribuição bootstrap do coeficiente de Spearman de cada achado, com 4.000 reamostras com reposição; a barra marca o intervalo de 95\% e o ponto, o valor da amostra. \textbf{b}, Coeficiente calculado separadamente nas duas metades de uma partição estrutural --- usinas cuja barra de conexão consta do cadastro de rede de 230 kV ou mais do ONS, e as demais --- com o valor $p$ do teste $z$ de Fisher para a diferença entre metades. \textbf{c}, O mesmo coeficiente calculado no grafo da EPE, usado nas Figs. 10 e 11, e no grafo do ONS, construído por número de barra; abaixo de cada barra, o tamanho da amostra no grafo alternativo, ou a marca de que o critério não se aplica --- a capacidade da usina não deriva do grafo e não tem versão alternativa. \textbf{d}, Placar: quantos dos critérios aplicáveis cada achado passa; o traço marca o critério não aplicável. Dados: Figs. 10 e 11; ONS, linhas de transmissão e subestações da Rede de Operação e fator de capacidade. Dados-fonte em \texttt{data/processed/fig12\_source\_robustez.csv}.}

\begin{document}

%================================================================ CAPA
\vspace*{6mm}
{\color{claro}\textsf{ANÁLISE DE DADOS ABERTOS DO SETOR ELÉTRICO BRASILEIRO}}\\[3mm]
{\fontsize{30}{34}\selectfont\textbf{O corte fotovoltaico no SIN}}\\[2mm]
{\fontsize{16}{20}\selectfont\color{cinza} Magnitude, composição da causa e alocação entre usinas, em doze análises sobre dados abertos}\\[6mm]
{\color{claro}\rule{\linewidth}{0.4pt}}\\[4mm]

\begin{minipage}[t]{0.47\linewidth}\vspace{0pt}
\bloco{Escopo}
O objeto é o corte de geração fotovoltaica no Sistema Interligado Nacional: quanto
se corta, por que motivo declarado e como o corte se distribui entre usinas.

O objetivo inicial --- detectar instalações fotovoltaicas por satélite para verificar
programas de crédito rural --- foi redefinido após duas medidas. A \textbf{Fig.\,1}
mede que o corte incide sobre toda a janela solar e é majoritariamente de razão
energética, de modo que a geração verificada não é proxy do recurso solar. A
\textbf{Fig.\,2} mede que 91\% da micro e minigeração distribuída é sub-pixel no
Sentinel-2, o que inviabiliza a verificação por imagem com sensor gratuito. A
\textbf{Fig.\,8} desloca a análise para a decisão de despacho.

\bloco{Revisões internas}
Cinco das doze análises examinam resultados das anteriores: duas revisam uma
conclusão (Figs.\,6 e 9) e três delimitam escopo (Figs.\,10, 11 e 12). As conclusões
revisadas permanecem no texto com a correção correspondente.

A \textbf{Fig.\,6} revisa a conclusão secundária da Fig.\,5: o viés atribuído à
geração de referência do ONS decorre do contrafactual adotado --- 2023 já era um ano
com restrição, anterior ao início da apuração pública.

A \textbf{Fig.\,9} retira a atribuição causal da Fig.\,8: decomposta por subsistema, a
coincidência entre corte e térmica inflexível não sustenta leitura local em nenhum dos
dois subsistemas com corte.

\end{minipage}\hfill
\begin{minipage}[t]{0.47\linewidth}\vspace{0pt}
\bloco{Conteúdo}
\footnotesize
\begin{tabular}{@{}p{11mm}p{58mm}p{28mm}@{}}
\textbf{Método} & Janela, amostra, teste e multiplicidade & \textit{referência} \\
\textbf{Fig.\,1} & Corte fotovoltaico no SIN & mecanismo \\
\textbf{Fig.\,2} & MMGD: detectabilidade por satélite & viabilidade \\
\textbf{Fig.\,3} & Estudo de caso pareado, PIE $\times$ APE & caso \\
\textbf{Fig.\,4} & Cronologia do entorno & hipótese temporal \\
\textbf{Fig.\,5} & Controle climático & restrição ou recurso \\
\textbf{Fig.\,6} & Calibração da referência do ONS & \textit{revisa a 5} \\
\textbf{Fig.\,7} & Detecção de construção por Sentinel-2 & executa $H_2$ \\
\textbf{Fig.\,8} & Despacho fora da ordem de mérito & mecanismo \\
\textbf{Fig.\,9} & Decomposição por subsistema & \textit{revisa a 8} \\
\textbf{Fig.\,10} & Topologia da rede de transmissão & \textit{responde à 3} \\
\textbf{Fig.\,11} & Replicação em eólica & \textit{delimita a 10} \\
\textbf{Fig.\,12} & Robustez dos resultados de rede & \textit{qualifica 10 e 11} \\
\textbf{Aux} & Imagem comercial: resolução e custo & decisão \\
\textbf{Aplic.} & Produção e exposição de uma área nova & exercício \\
\end{tabular}
\normalsize

\bloco{Fontes}
Exclusivamente dados abertos, sem autenticação: ONS (\textit{constrained-off},
geração térmica por motivo de despacho, fator de capacidade, curva de carga,
intercâmbio, energia armazenada, CMO), ANEEL (SIGA e geração distribuída),
Sentinel-2 L2A pelo catálogo STAC público, NASA POWER e malhas do IBGE.

\end{minipage}

\clearpage

%================================================================ MÉTODO
\cab{MÉTODO}{Janela, amostra e multiplicidade}{O recorte de cada figura, a unidade que se repete dentro dele, o teste que ela declara, o contrato de exportação e o número total de testes reportados.}

\begin{minipage}[t]{0.635\linewidth}\vspace{0pt}
\bloco{Janela, amostra, unidade de repetição e teste}
\par\vspace{3pt}\small
\begin{tabular}{@{}lp{41mm}p{26mm}p{31mm}p{49mm}@{}}
\toprule
 & janela & $n$ & unidade de repetição & teste declarado \\
\midrule
\textbf{1} & abr/2024--jul/2026 & 85 usinas & usina · 30 min & nenhum: descritiva \\
\textbf{2} & conexões até abr/2026 & 4,6 M instalações & empreendimento & nenhum: descritiva \\
\textbf{3} & out/2024--ago/2026 (par) & 2 ativos; 60 conj. & conjunto & Spearman; 5 covariáveis, sem correção \\
\textbf{4} & 2023--jul/2026; entorno 2021 & 43 e 24 meses & conjunto · mês & nenhum teste formal \\
\textbf{5} & 2023--jul/2026 & meses com clima & conjunto · mês & nenhum: decomposição aritmética \\
\textbf{6} & abr/2024--jul/2026 & 74 calibradas & usina · dia-usina & Spearman; 2 covariáveis \\
\textbf{7} & 2018--2026 & 34 conj., 306 med. & conjunto · ano & Wilcoxon pareado; 2 comparações \\
\textbf{8} & abr/2024--ago/2026 & 21.192 horas & hora do SIN & Mann-Whitney e Spearman; quintis \\
\textbf{9} & idem, 9--15\,h & 21.192 horas & hora do subsistema & Spearman; sem correção \\
\textbf{10} & set/2025--ago/2026 & 59 conjuntos & conjunto & Spearman; 8 covariáveis, \textbf{Holm} \\
\textbf{11} & nov/2022--jul/2025 eól. & 120 conjuntos & conjunto & idem, \textbf{Holm} por fonte \\
 & set/2025--ago/2026 sol. & 59 conjuntos & conjunto &  \\
\textbf{12} & amostras das Figs.\,10 e 11 & 59 e 120 & conjunto & bootstrap e $z$ de Fisher \\
\bottomrule
\end{tabular}
\par\vspace{3pt}
\ressalva{As janelas de corte saem da mesma regra --- percentil 80 do início e 20 do fim da apuração --- e por isso diferem entre fontes. A Fig.\,1 vem de extração anterior às demais. Cada figura repete estas declarações na sua própria página, com as exclusões.}
\end{minipage}\hfill
\begin{minipage}[t]{0.335\linewidth}\vspace{0pt}\small
\bloco{Padrão de figura}
Todas as figuras seguem o mesmo contrato: 183\,mm de largura, corpo 7\,pt, piso de
5\,pt para qualquer glifo, texto editável em PDF e SVG, TIFF a 600\,dpi, dados-fonte
exportados em CSV e auditoria automática de colisão obrigatória. Cada figura declara
$n$, unidade de repetição, teste, exclusões e limites de interpretação.

\bloco{Multiplicidade}
A correção é aplicada \textbf{dentro} de figura --- Holm sobre as oito covariáveis
pré-declaradas nas Figs.\,10 e 11 --- e declarada ausente na Fig.\,3, onde as cinco
covariáveis testadas resultaram todas nulas. \textbf{Não há correção entre figuras.}
Somados os testes que vão a arquivo, a série reporta \textbf{49} --- 5 na Fig.\,3, 2 na
7, 5 na 8, 14 na 9, 8 na 10, 8 na 11 e 7 na 12 ---, dos quais \textbf{16 passam por Holm
e 33 não}. As oito covariáveis foram pré-declaradas e testadas na mesma ordem nas duas
fontes, e a Fig.\,12 submete as sobreviventes a testes adicionais. Um limiar de
multiplicidade em nível de série não é aplicado neste documento.
\end{minipage}
\clearpage

%================================================================ FIG 1
\analise{ANÁLISE 1 \quad · \quad ONS, restrição por constrained-off}
{O corte fotovoltaico no SIN é sistêmico}
{Incide sobre toda a janela solar, cresce em ritmo semelhante no Nordeste e no Sudeste/CO, e é majoritariamente de razão energética, classificada como origem sistêmica.}
{../figures/fig1_curtailment_sin.pdf}
{
\bloco{Método}
O corte é definido como
\[ \max\bigl(\text{referência} - \text{verificada},\; 0\bigr) \]
nos intervalos com razão de restrição apurada. O campo \texttt{val\_geracaolimitada}
\textbf{não} é a energia cortada: é o teto ao qual a geração foi limitada. Sua
correlação com a geração verificada é 0,86 e com a perda apurada é \textbf{0,02}.
Usá-lo como energia cortada superestima a perda.

\bloco{Resultados}
Corte de \textbf{21,6\%} no parque com apuração. A razão energética domina o total
cortado; a origem declarada é predominantemente sistêmica.

\bloco{Interpretação}
Onde a geração de uma usina é limitada por decisão operativa, um modelo que preveja
geração a partir de irradiância ou de imagem de satélite ajusta um alvo que incorpora
essa decisão, e não apenas o recurso físico. Isso condiciona o alvo de modelagem de
todas as análises seguintes.

\lit{Os 21,6\% medidos são de solar centralizada com apuração, na janela
abr/2024--jul/2026. O ONS reporta 20,7\% para eólica \textit{e} solar em 2025 --- ordens
de grandeza compatíveis, denominadores diferentes. Os valores publicados para outros
sistemas são menores: o Chile perdeu 17\% da renovável variável disponível em 2024, a
Espanha registrou $\approx$11\% em jul/2025, e na Califórnia a curtailment \textit{média}
da solar é 4,3\%. Novan e Wang (2024) medem que a taxa \textbf{marginal} é cerca de duas
vezes a média, o que implica exposição de uma usina nova superior a 21,6\% caso a razão
se mantenha no Brasil.}

\bloco{Limitações}
\ressalva{O denominador é o das usinas com apuração, não do parque solar nacional.}
\ressalva{Não cobre MMGD, que tem regime de restrição distinto.}
\ressalva{A referência é estimativa do operador: em 21,8\% dos intervalos restritos ela fica abaixo da geração verificada, e o truncamento em zero torna a estimativa conservadora, não exata.}
\ressalva{Razão e origem são classificações operativas do ONS, não inferência causal independente.}
}{\legI}

%================================================================ FIG 2
\analisec{ANÁLISE 2 \quad · \quad ANEEL, geração distribuída}
{91\% da MMGD é sub-pixel no Sentinel-2}
{A verificação de instalações registradas por imagem de satélite é inexequível com sensor gratuito, e a data de conexão não é exógena.}
{../figures/fig2_mmgd_detectabilidade.pdf}
{
\bloco{Método}
A ANEEL publica a \textbf{área do arranjo} de cada instalação. Sobre 4,6 milhões de
registros, o número de pixels que cada sensor entrega por instalação é
\[ \mathrm{px}(s) = A_{\text{arranjo}} / r_s^2 \]

\bloco{Resultados}
Mediana de pixels por instalação e fração sub-pixel: Sentinel-2 (10\,m)
\textbf{0,29\,px} e \textbf{91,2\%}; CBERS-4A WPM (2\,m) 7,25\,px e 0,5\%; VHR
comercial (0,5\,m) 116\,px e $\approx$0. A classe rural --- alvo dos programas de
crédito --- tem área mediana de 41\,m$^2$ e \textbf{5,2\%} de segmentabilidade no
Sentinel-2.

\bloco{Interpretação}
\textbf{Localização.} O registro \textbf{não tem localização utilizável}: os campos de
coordenada existem no esquema e estão preenchidos em \textbf{0,06\%}. Há
\texttt{CodCEP} --- completo para os 7,7\% de pessoa jurídica, truncado em 5 dígitos
para os 92,3\% de pessoa física ---, com mediana de 1 prefixo por município, o que
\textbf{não refina o alvo rural}. Sem coordenada por instalação, a verificação exige
varredura exaustiva; apenas o segmento de pessoa jurídica é endereçável. A resolução do
sensor não é a restrição dominante.

\textbf{Exogeneidade da data de conexão.} A janela de transição da Lei 14.300/2022
concentrou 98 mil conexões em um único mês --- quatro vezes a média do ano anterior. A
data de conexão não é exógena: um estudo de evento sobre ela mediria a antecipação
regulatória, não a política.

\lit{A literatura que segmenta telhados opera em 0,1--0,3\,m: DeepSolar (Yu et al., 2018)
usa imagem aérea submétrica; o conjunto multirresolução de Jiang et al.\ (2021) vai de
0,1 a 0,8\,m. O inventário global de Kruitwagen et al.\ (2021), em \textit{Nature}, é a
referência com Sentinel-2 e cobre apenas instalações acima de 10\,kW. Não foi localizado
trabalho que segmente telhado residencial a 10\,m. A Fig.\,2 quantifica a razão.}

\bloco{Limitações}
\ressalva{Contar pixels não é detectar: um objeto sub-pixel ainda altera a assinatura espectral. O que a contagem descarta é a segmentação.}
\ressalva{O piso de 4 pixels é convenção declarada, não constante física.}
\ressalva{É censo administrativo, não amostra: não há erro amostral a declarar.}
}{\legII}

%================================================================ FIG 3
\analisec{ANÁLISE 3 \quad · \quad ONS + ANEEL/SIGA, estudo de caso pareado}
{Diferença de 3,0$\times$ entre dois ativos, sem covariável local que a reproduza na frota}
{Sol do Cerrado perde 3,0$\times$ mais energia que Três Marias 3. Nenhuma das cinco covariáveis locais testadas reproduz esse padrão no parque.}
{../figures/fig3_estudo_caso.pdf}
{
\bloco{Método}
Sol do Cerrado (681,3\,MW, produtor independente, 17 SPEs de um grupo, barramento
de Jaíba 230\,kV \textit{dividido}) contra Conj.\ Três Marias 3 (70,0\,MW,
autoprodução, sete empresas privadas distintas, ponto \textit{exclusivo}). Ambos em
Minas Gerais, mesmo subsistema, janela comum recortada programaticamente.

O campo \texttt{nom\_agenteproprietario} do ONS identifica o \textbf{agente
representante}, não o proprietário, e não sustenta a leitura de que Três Marias 3
seja ativo estatal. O SIGA registra sete proprietários privados sob regime de
autoprodução; não há usina solar estatal no parque despachado, e o eixo público
$\times$ privado não é testável nesta amostra.

\bloco{Resultados}
Taxa de corte de 26,8\% contra 8,8\%. A energia cortada equivale a cerca de
\textbf{R\$\,64 milhões} contra pouco mais de R\$\,1 milhão, valorada a CMO.

Nos dois ativos o CMO ponderado \textbf{durante o corte} é menor que durante a
geração, o que é consistente com restrição de razão energética. A Fig.\,8 examina
esse ponto no despacho.

\bloco{Interpretação}
Duas explicações locais eram plausíveis: barramento compartilhado e densidade
regional de renovável. Nenhuma se reproduz na frota --- nem carga do ponto de conexão,
nem número de conjuntos, nem densidade regional, nem tensão, nem capacidade própria.
Uma explicação local consistente com um par de ativos não se generaliza a mecanismo
sem esse teste.

\bloco{Revisão}
A versão anterior desta análise concluía que a topologia de transmissão não está em
dado aberto. \textbf{Está.} O SIGET publica o grafo linha a linha --- 1.166 linhas,
619 subestações com código ONS, tensão e comprimento --- e o \texttt{fator-capaci\-dade-2}
traz a coordenada do ponto de conexão de cada usina; o join cobre 168 dos 229
conjuntos despachados. O \textbf{limite de transferência} também está publicado:
\texttt{linha-transmissao} do ONS traz capacidade operativa em MVA com variantes
sazonais, além de resistência e reatância. Permanece indisponível
\textbf{qual restrição está ativa em cada intervalo}. As cinco covariáveis testadas
são locais ao ponto e não usam a estrutura da rede, de modo que o resultado nulo não
se estende a ela; a \textbf{Fig.\,10} a testa.

\lit{Newbery (2024): quem é cortado depende do regime de acesso, não apenas da física da rede.}

\bloco{Limitações}
\ressalva{Dois ativos não são amostra: o par foi escolhido por contraste.}
\ressalva{Cinco covariáveis testadas sem correção de multiplicidade; como nenhuma resultou significativa, a correção só tornaria a conclusão mais conservadora.}
\ressalva{CMO é proxy de valor sistêmico, não preço de liquidação nem prejuízo líquido.}
}{\legIII}

%================================================================ FIG 4
\analise{ANÁLISE 4 \quad · \quad ONS, fator de capacidade}
{Cronologia do corte e da energização do entorno}
{A apuração de corte começa em abr/2024 e a onda de construção em Jaíba é de 2021 a 2024. O fator de capacidade cobre o período anterior à apuração.}
{../figures/fig4_cronologia.pdf}
{
\bloco{Método}
A série de corte começa em abr/2024 e não cobre o período em que a hipótese se
manifestaria. O fator de capacidade existe desde a entrada em operação e cai quando a
geração é limitada:
\[ FC_{\text{verificado}}(t)=\frac{E_{\text{gerada}}(t)}{P_{\text{inst}}\cdot\Delta t} \]
A partir de abr/2024 a apuração permite recompor o potencial:
\[ FC_{\text{potencial}} = FC_{\text{verificado}} + FC_{\text{cortado}} \]
Se a queda for de restrição de rede, o potencial recomposto permanece estável
enquanto o verificado cai. Se for recurso solar ou degradação, os dois caem juntos.

\bloco{Interpretação}
A decomposição separa restrição de degradação, mas não separa restrição de
\textbf{variação climática}: se 2025 tiver sido menos ensolarado que 2023, a queda
não informa sobre a rede. A Fig.\,5 mede a irradiância diretamente sobre cada usina.

\lit{O desenho é um estudo de evento com adoção escalonada --- cada vizinho energiza em
data diferente ---, caso em que Callaway e Sant'Anna (2021) mostram que o estimador de
diferenças em diferenças com efeitos fixos duplos fica enviesado, por ponderar
comparações entre unidades já tratadas. A leitura desta figura é descritiva e pareada
contra um único controle, o que evita o problema e também não estima efeito de
tratamento. Estimá-lo exigiria o estimador por grupo e período daqueles autores e uma
definição de tratamento exógena, que a Fig.\,2 mostra não existir do lado da MMGD.}

\bloco{Limitações}
\ressalva{O fator de capacidade é observável de longo prazo, mas responde a várias causas ao mesmo tempo; a decomposição só vale onde a apuração existe.}
\ressalva{Comparação pareada sem aleatorização, como na Fig.\,3.}
}{\legIV}

%================================================================ FIG 5
\analise{ANÁLISE 5 \quad · \quad ONS + NASA POWER}
{A queda do fator de capacidade não é atribuível à irradiância}
{Normalizada pela irradiância medida sobre a usina, a perda de Sol do Cerrado permanece: nove décimos da queda não são atribuíveis ao recurso solar.}
{../figures/fig5_controle_climatico.pdf}
{
\bloco{Método}
Com a irradiância global horizontal diária $H$ sobre a usina, a razão de desempenho
\[ PR = \frac{FC \cdot 24\,\mathrm{h}}{H} \]
divide a energia gerada por kW instalado pelo recurso disponível. Se a queda fosse de
menor irradiância, $PR$ permaneceria estável.

\bloco{Resultados}
A irradiância sobre Jaíba caiu \textbf{2,8\%} entre 2023 e 2025. O fator de
capacidade caiu \textbf{29,1\%}. Escalando o FC de 2023 pela razão de irradiância, o
esperado para 2025 seria 0,206; o observado foi 0,150. Cerca de um décimo da queda é
atribuível ao recurso.

A razão de desempenho cai de 0,854 para 0,619 em Sol do Cerrado. No controle, Três
Marias 3 permanece em torno de 0,99 sem tendência.

\bloco{Interpretação}
Os meses de 2025 não estão na curva de 2023 em posição de menor irradiância: estão em
uma curva mais baixa. O deslocamento é consistente com perda independente do recurso.

\bloco{Revisão}
O déficit não explicado pela irradiância é de 0,056 em fator de capacidade; o corte
apurado pelo ONS equivale a 0,147 --- 2,6 vezes maior. Essa diferença foi interpretada
como indício de que a referência do ONS superestima o contrafactual. A interpretação
pressupõe que 2023 fosse um ano sem restrição; a Fig.\,6 mostra que não era, e a
conclusão secundária não se sustenta.

\lit{A irradiância entra como denominador da razão de desempenho, na condição de variável
derivada de sensoriamento. Proctor, Carleton e Sum (2023) mostram que usar uma variável
assim em inferência a jusante enviesa o ponto e o erro padrão mesmo com erro de medida
pequeno, e recomendam tratá-la como distribuição --- imputação múltipla --- e não como
valor pontual. Alix-García e Millimet (2023) acrescentam que o erro de produto de satélite
raramente é clássico. Nenhuma das duas correções foi aplicada: a diferença medida é grande
o bastante para sobreviver a erro moderado, mas o intervalo em torno dela é mais estreito
do que deveria ser.}

\bloco{Limitações}
\ressalva{Reanálise, não piranômetro: NASA POWER é MERRA-2 combinado com satélite, em célula de 0,5$^\circ\times$0,625$^\circ$.}
\ressalva{GHI no lugar de POA: a irradiância é horizontal, os módulos não. O PR aqui é indicador relativo dentro de cada usina, não PR de norma IEC.}
\ressalva{Temperatura não modelada; contrafactual linear.}
}{\legV}

%================================================================ FIG 6
\analise{ANÁLISE 6 \quad · \quad revisão da Fig.\,5}
{A geração de referência do ONS é calibrada}
{Em 74 usinas o viés mediano é 1,02. O déficit aparente da Fig.\,5 decorre do contrafactual, não do operador: 2023 já era um ano com restrição.}
{../figures/fig6_vies_referencia.pdf}
{
\bloco{Método}
Para cada usina, calibra-se uma razão de desempenho \textbf{nos dias sem restrição da
própria usina} ($f<1\%$) e compara-se o potencial implicado pelo ONS (verificado $+$
cortado) com o potencial previsto pela irradiância local.

\bloco{Resultados}
Viés mediano por usina \textbf{1,02}, com metade entre 0,98 e 1,06. Nenhuma das 74
usinas calibradas passa de 1,5, e 26 ficam abaixo de 1 --- para um terço delas a
referência do ONS é \textit{conservadora}. No agregado energético dos dias restritos,
o corte apurado é 1,03 vezes o implicado pela irradiância. A associação entre viés e
intensidade de corte é nula, o que descarta a falha de maior consequência.

\bloco{Revisão}
O desempenho limpo de Sol do Cerrado, calibrado em 279 dias sem restrição, é
\textbf{0,931}. Em 2023 a usina operou a 0,854 --- \textbf{8,3\% abaixo do seu próprio
desempenho limpo}. Aquele ano já era restrito, sem apuração pública. Ao usar 2023 como
base, a Fig.\,5 mediu um déficit reduzido e concluiu, incorretamente, que o ONS
superestimava o corte.

\bloco{Interpretação}
A conclusão principal da Fig.\,5 --- perda de rede, não de recurso --- permanece e é
reforçada: contra o desempenho limpo, o corte implícito é \textbf{33\% em 2025}, não
26\%. A afirmação secundária sobre viés do operador é retirada. Os 21,6\% da Fig.\,1
não requerem correção: com a referência calibrada, o agregado permanece válido.

\lit{A geração de referência do ONS é derivada de curvas de desempenho da usina
(irradiância ou vento contra potência), como descrevem Vieira et al.\ (2026). Validá-la
por fora exige fonte independente: no lado eólico, Silva et al.\ (2026) usam anemometria
de SCADA. A verificação equivalente no lado solar, com irradiância de satélite, não foi
localizada na literatura consultada; a busca não foi sistemática.}

\bloco{Limitações}
\ressalva{Dias ``limpos'' podem não ser limpos; como o resultado é ausência de viés, esse erro só o tornaria mais forte.}
\ressalva{Seis usinas ficaram de fora por não atingirem 30 dias limpos --- justamente as mais restritas, onde a referência é mais usada.}
\ressalva{O $PR$ limpo é constante por usina e herda qualquer viés sazonal dos dias de calibração.}
}{\legVI}

%================================================================ FIG 7
\analisec{ANÁLISE 7 \quad · \quad Sentinel-2 L2A (STAC público) + ANEEL/SIGA}
{A construção é detectável por Sentinel-2; a operação não é}
{A construção de uma usina fotovoltaica eleva a razão de brilho no vermelho em 34\% na mediana e data a entrada em operação; a fase operacional não se separa do entorno a 10\,m.}
{../figures/fig7_deteccao_construcao.pdf}
{
\bloco{Método}
Para cada conjunto, mede-se anualmente a razão entre o brilho no vermelho dentro do
arranjo e num anel de controle:
\[ \rho(t)=\frac{\overline{B_{\text{red}}}\big|_{d\le 400\,\mathrm{m}}}
{\overline{B_{\text{red}}}\big|_{1{,}5\le d\le 3\,\mathrm{km}}} \]
O anel controla iluminação, estação e umidade do solo. As séries são alinhadas pela
\textbf{data de entrada em operação} do SIGA, não pelo calendário. 34 conjuntos,
9 anos, 306 medições.

\bloco{Resultados}
Da base à construção a razão sobe de 1,01 para 1,30 --- \textbf{+34\% na mediana, 24 de
30 usinas}, Wilcoxon $p=1{,}2\times10^{-5}$. Da base à operação a mediana vai de 0,95
para 1,02 --- \textbf{$+6{,}6\%$}, com 11 de 18 usinas em alta contra 7 em queda
($p=0{,}09$).

\bloco{Interpretação}
A construção remove vegetação de centenas de hectares, e solo exposto tem refletância
mais alta no vermelho. Na fase operacional não há sinal em nenhuma direção: módulos são
escuros, mas a 10\,m o pixel mistura painel, corredor, via e solo.

A hipótese ``registro $\leftrightarrow$ evidência no satélite'' é executável para
geração centralizada, com uma inversão: o satélite confirma a usina pela cicatriz da
obra, não pelo arranjo. Trata-se de evento transitório, de modo que a auditoria exige
monitoramento contínuo, não observação do estado atual.

\bloco{Revisão}
A primeira execução usou a coordenada do \textbf{conjunto} publicada pelo ONS, que
corresponde à subestação coletora, e resultou nula. Em \textbf{31 de 37} conjuntos a
subestação fica a mais de 750\,m do centroide do arranjo, com mediana de 1,9\,km, de
modo que a máscara amostrava terreno fora do arranjo. Refeita com as coordenadas de
usina do SIGA, a diferença entre fases aparece.

\lit{A razão dentro/fora do arranjo é uma versão simplificada do \textit{Temporal Cluster
Matching} (Robinson et al., Microsoft Research), que data construções comparando a relação
espectral dentro e fora da pegada, e foi aplicado a usinas solares em Karnataka. O Global
Renewables Watch estima datas de construção em escala global por imagem de alta resolução.
A Fig.\,7 acrescenta a validação contra o \textbf{registro de outorga} e o resultado
negativo sobre a fase operacional a 10\,m.}

\bloco{Limitações}
\ressalva{A qualidade da geolocalização de referência domina o resultado; qualquer replicação depende dessa camada.}
\ressalva{Uma cena por ano fixa a datação em resolução anual: acerto exato em 21\% e dentro de um ano em 62\% dos casos.}
\ressalva{Círculos de 400\,m não são o polígono do arranjo; delinear por segmentação melhoraria o contraste.}
\ressalva{O pulso indica remoção de vegetação, que também ocorre em lavoura, mineração e loteamento: a atribuição depende do cruzamento com o registro. O resultado não se estende a telhados.}
}{\legVII}

%================================================================ FIG 8
\analise{ANÁLISE 8 \quad · \quad ONS, geração térmica por motivo de despacho}
{O corte por excesso de oferta coincide com térmica inflexível}
{Dois terços do corte são classificados como excesso de oferta e ocorrem com 4,6\,GWmed de térmica fora da ordem de mérito no mesmo instante, 73\% dela inflexibilidade declarada.}
{../figures/fig8_despacho_merito.pdf}
{
\bloco{Método}
A decomposição do ONS é \textbf{aninhada, não aditiva}: somar todos os campos dá
1,58 vez a geração total. A definição consistente é
\[ \text{fora do mérito} = \texttt{verifgeracao} - \texttt{verifordemmerito} \]
verificada contra a soma dos motivos não-econômicos disjuntos --- razão
\textbf{1,0030}, mediana da diferença 0,17\,MWmed.

\bloco{Resultados}
21.889\,GWh de corte na janela, 68,5\% por razão energética. Nas horas com esse corte
há 4.564\,MWmed de térmica fora do mérito contra 3.146 nas demais (+45\%); ponderado
pela energia cortada, 5.266. \textbf{73\% é inflexibilidade pura.}

Da madrugada ao pico solar a parcela por mérito cai \textbf{1.640\,MWmed} e a parcela
fora do mérito sobe \textbf{1.184\,MWmed}: a mesma térmica permanece ligada e deixa de
ser econômica quando o CMO cai.

A correlação entre corte e térmica inflexível é $\rho=0{,}48$ no quintil de carga mais
baixa, 0,38 no segundo, e 0,09--0,20 nos três superiores. Com $n\approx1.240$ por
estrato \textbf{todos os quintis dão $p<0{,}05$}; a leitura se apoia na magnitude, não
na significância, e a queda não é monotônica.

\bloco{Interpretação}
A leitura de que a térmica desloca a geração solar não sobrevive ao controle de carga:
a carga sobe aproximadamente o quanto a geração solar sobe, e a \textbf{carga líquida é
plana} (72,2 $\rightarrow$ 72,1\,GWmed). O que varia é a classificação do despacho
térmico, não o seu volume.

\lit{O mecanismo é documentado na literatura chinesa de curtailment: unidades térmicas com
mínimo técnico e status \textit{must-run} ocupam a faixa que a renovável variável
utilizaria, e a resposta de política foi o programa de \textit{flexibility retrofit} ---
16,35\,GW em demonstração em 2017, somando 4,8\,GW de capacidade de regulação de ponta.
No Brasil, Vieira et al.\ (2026) modelaram o corte com regressão logística horária e
identificaram a MMGD como preditor mais forte do corte por demanda ($OR\approx9{,}8$), com
a carga do sistema atuando de forma protetiva --- o gradiente que o painel \textbf{c}
reproduz.}

\bloco{Limitações}
\ressalva{Coincidência não é causalidade: separar exigiria reotimizar o despacho sem a inflexibilidade declarada.}
\ressalva{Inflexibilidade é declaração do agente, não mínimo técnico verificável em dado aberto.}
\ressalva{Agregado nacional --- revisado na Fig.\,9.}
\ressalva{Controle negativo: energia armazenada não explica ($\rho=0{,}06$, $p=0{,}06$).}
}{\legVIII}

%================================================================ FIG 9
\analise{ANÁLISE 9 \quad · \quad revisão da Fig.\,8}
{O piso térmico inflexível não coincide geograficamente com o corte}
{Decomposta por subsistema, a coincidência nacional não sustenta leitura causal local em nenhum dos dois subsistemas com corte.}
{../figures/fig9_decomposicao_subsistema.pdf}
{
\bloco{Resultados}
Nordeste e Sudeste cortam volumes semelhantes nas horas solares; a térmica inflexível
não acompanha.
\par\vspace{1pt}\scriptsize
\begin{tabular}{@{}lrr@{}}
\toprule
MWmed, horas de 9\,h a 15\,h & Nordeste & Sudeste/CO \\
\midrule
corte por razão energética & 1.079 & 950 \\
inflexibilidade pura local & \textbf{36} & 1.739 \\
corte / piso local & \textbf{30,2$\times$} & 0,5$\times$ \\
$\rho$ com térmica do NE & +0,08 & +0,10 \\
$\rho$ com térmica do SE & \textbf{+0,42} & \textbf{+0,41} \\
\bottomrule
\end{tabular}
\par\vspace{2pt}
No Nordeste, onde o volume de corte é o maior do país, há 36\,MWmed de inflexibilidade
contra 1.079 de corte; a associação local é $\rho=0{,}08$ e não passa de 0,16 em nenhum
quintil de carga.

No Sudeste/CO o painel \textbf{b} reproduz o padrão da Fig.\,8, com $\rho=0{,}40$ nos
quintis de carga baixa. O painel \textbf{c} não sustenta a leitura local: a térmica do
Sudeste se correlaciona com o corte do Sudeste em 0,41 e com o do \textbf{Nordeste} em
0,42, a mais de dois mil quilômetros. O valor de $\rho$ é determinado por qual térmica
se mede, não por qual corte; uma causa local produziria associação mais forte com o
corte local.

\bloco{Interpretação}
A hipótese de saturação do corredor também não se sustenta. Nas horas de corte alto o
fluxo NE$\rightarrow$SE é \textit{menor} --- 2.993 contra 3.848\,MWmed --- e fica acima
do teto empírico em 2,9\% das horas contra 5,7\% nas demais ($\rho=-0{,}25$): o
Nordeste é cortado com o corredor com folga.

\bloco{Revisão}
Nenhum valor da Fig.\,8 muda. O que se retira é a \textbf{atribuição causal}: a
coincidência é real e compatível com um fator comum --- carga nacional baixa --- e este
desenho não distingue as duas hipóteses.

\lit{Newbery (2024) trata de \textit{quem} é cortado sob restrição de transmissão e de como
o regime de acesso distorce o sinal de investimento, o que fornece o enquadramento
econômico do descompasso medido. A invisibilidade da MMGD ao despacho é problema
documentado: geração atrás do medidor não é observada pelo operador e eleva o erro de
previsão de carga líquida em 30\% no nível domiciliar e até 250\% no agregado.}

\bloco{Limitações}
\ressalva{Subsistema ainda é agregação grosseira: o Sudeste/CO vai de Minas ao Mato Grosso do Sul.}
\ressalva{O teto do corredor é o percentil 95 do próprio fluxo. Verificado: \texttt{programacao\_fluxo\_controlado} cobre só HVDC e interligações internacionais, e \texttt{linha-transmissao} dá capacidade por linha, não do corredor --- que resulta de estudo de estabilidade.}
\ressalva{``Ausência de demanda'' é inferência por eliminação, não testada positivamente.}
\ressalva{O painel c não estabelece ausência de mecanismo local no Sudeste: mostra que este desenho não o identifica, porque inflexibilidade local e carga nacional são colineares.}
}{\legIX}

%================================================================ FIG 10
\analisea{ANÁLISE 10 \quad · \quad EPE (rede) + ONS \quad · \quad responde à Fig.\,3}
{A distância de rede até a demanda está associada à taxa de corte}
{Com o grafo de transmissão montado, duas de oito covariáveis pré-declaradas sobrevivem à correção de multiplicidade.}
{../figures/fig10_topologia_rede.pdf}
{
\bloco{Método}
A Fig.\,3 testou cinco covariáveis \textbf{locais ao ponto de conexão} e nenhuma
resultou significativa. As oito covariáveis testadas aqui são \textbf{estruturais}:
dependem do grafo. Rede da EPE: 920 subestações, 1.896 linhas resolvidas em aresta
(87\%), ancoragem mediana de 0,10\,km. Holm sobre as oito hipóteses pré-declaradas.

\bloco{Resultados}
Duas sobrevivem a Holm, com magnitude próxima e sinais opostos: \textbf{distância de
rede à carga} ($+$0,38) e \textbf{capacidade da usina} ($-$0,38), ambas com Holm
$p = 0{,}023$. A soma de tensões no nó atinge $p < 0{,}05$ bruto e não sobrevive à
correção; as outras cinco não atingem.

A distância mediana à carga é próxima nos dois subsistemas --- 497\,km no Nordeste,
520 no Sudeste/CO --- e a associação se mantém \textbf{dentro} do Nordeste
($\rho = +0{,}34$, $p = 0{,}027$, $n = 43$), mas não dentro do Sudeste/CO
($+0{,}31$, $p = 0{,}25$, $n = 16$). A capacidade da usina se mantém nos \textit{dois}
--- $-0{,}41$ e $-0{,}44$ ---, no Nordeste com $p = 0{,}006$. A estrutura local do nó
não apresenta associação.

\bloco{Interpretação}
Quanto maior a distância em quilômetros de linha até o centro de carga, maior a taxa de
corte. O resultado é consistente com a Fig.\,9, em que o Nordeste é cortado com o
corredor NE$\rightarrow$SE com folga: as duas medidas apontam para a distância até a
demanda, não para a saturação local do nó.

A capacidade da usina tem sinal \textbf{negativo}: as maiores são cortadas menos. Não
era hipótese direcional. As explicações plausíveis --- qualidade de conexão, seleção
contratual --- são testáveis e \textbf{não foram testadas}.

\bloco{Escopo e robustez}
O resultado vale para fotovoltaica: a replicação em eólica não se sustenta, e a
Fig.\,11 mede a razão. A Fig.\,12 avalia a robustez das duas covariáveis: a distância
à carga tem intervalo bootstrap de largura 0,49, coeficiente que cai de 0,62 para 0,11
conforme a metade da amostra, e não aparece no grafo construído por barra; a capacidade
da usina passa nos dois critérios que se aplicam a ela. Os três testes usam apenas
dados já disponíveis nesta análise.

\lit{Maji et al.\ (e-Energy 2025) sustentam que o corte é fenômeno de \textbf{nó} e que
estimá-lo exige descer ao ponto de conexão. É a hipótese que esta figura testa e que, no
fotovoltaico, não confirma: o preditor é a distância até a carga, grandeza de caminho, não
de nó. A Fig.\,11 mostra que em eólica a relação se inverte, o que delimita a afirmação
daqueles autores por composição da causa. Newbery (2024) fornece o enquadramento: se a
posição na rede determina a exposição ao corte e o regime de acesso não a precifica, o
sinal de investimento aponta para o lugar errado.}

\bloco{Limitações}
\ressalva{Distância de rede não é elétrica: caminho ponderado por km de linha. A impedância está em dado aberto --- 100\% das linhas de \texttt{linha-transmissao} (ONS), com capacidade em MVA em 86,7\% ---, e a escolha da camada geométrica da EPE impediu o uso da distância elétrica nesta figura.}
\ressalva{Centros de carga são nove metrópoles, proxy geográfico declarado.}
\ressalva{n = 59 para oito hipóteses. Holm foi aplicado, mas não há poder para efeitos pequenos. É associação, não causalidade.}
}{\legX}

%================================================================ FIG 11
\analisec{ANÁLISE 11 \quad · \quad ONS (eólica) + EPE \quad · \quad delimita a Fig.\,10}
{O resultado da Fig.\,10 não se replica em eólica}
{A covariável que prevê o corte de cada fonte acompanha a origem declarada do corte dessa fonte.}
{../figures/fig11_replicacao_eolica.pdf}
{
\bloco{Método}
O ONS apura \textit{constrained-off} eólico desde \textbf{out/2021} em \textbf{199
conjuntos}, contra 87 do solar. Mesmas oito covariáveis da Fig.\,10, mesma ordem, mesmo
grafo, Holm reaplicado dentro de cada fonte.

\bloco{Resultados}
A distância de rede à carga, que no solar dá $\rho = +0{,}38$ e sobrevive a Holm, dá
$\rho = -0{,}07$ em 120 conjuntos eólicos. A amostra eólica é o dobro da solar.

As covariáveis que sobrevivem em eólica são as que não sobreviveram em solar: usinas
no mesmo nó vai de $+0{,}07$ para \textbf{$+0{,}47$} e grau do nó de $-0{,}17$ para
\textbf{$+0{,}35$}, ambas sobrevivendo a Holm.

A classificação de origem do ONS, independente de todas as covariáveis testadas,
registra corte solar \textbf{91\% sistêmico} e eólico 76\%, com \textbf{24\% de origem
local} --- 2,7 vezes a fração do solar. Confiabilidade local é 7\% do corte solar e
\textbf{22\%} do eólico.

A amostra eólica é 91\% nordestina e a solar não é. Restringindo as duas ao Nordeste,
o solar mantém distância à carga em $+0{,}34$ e usinas no nó em $+0{,}05$; a eólica
apresenta $-0{,}19$ e $+0{,}36$. A dissociação se mantém sob controle de geografia.

\bloco{Interpretação}
A covariável que prevê o corte de cada fonte acompanha a origem declarada do corte
dessa fonte: corte de origem sistêmica é previsto por grandeza sistêmica --- distância
até a demanda ---, e corte de origem local, por grandeza local --- número de usinas
disputando o ponto.

Em termos operativos, a exposição solar responde a aproximação da demanda ---
armazenamento, deslocamento de carga --- e a eólica, a reforço local de escoamento.

\bloco{Revisão}
A Fig.\,10 enunciou a hipótese como sobre posição na rede, independentemente da
tecnologia. A replicação restringe o enunciado: vale para fonte cujo corte é
majoritariamente sistêmico. O resultado da Fig.\,10 permanece no fotovoltaico; muda a
generalidade.

Sob os testes da Fig.\,12, \textbf{grau do nó} passa nos três critérios e
\textbf{usinas no nó} replica no grafo alternativo, mas varia entre subamostras: a
direção se mantém, a magnitude não.

\lit{É o eixo canônico da literatura: Bird, Lew e Milligan (2016) organizam a revisão de
onze países por excesso de oferta sistêmico contra restrição local. A EirGrid formaliza a
distinção no vocabulário --- \textit{curtailment} para razão sistêmica, \textit{constraint}
para local --- e reporta 6,6\% contra 4,7\% no \textit{dispatch-down} eólico de 2025. Maji
et al.\ (e-Energy 2025) sustentam que o corte é fenômeno de nó. Não foi localizado teste
publicado de se a \textbf{covariável preditiva do corte varia com a composição da causa};
a busca não foi sistemática.}

\bloco{Limitações}
\ressalva{As janelas não coincidem --- eólica de nov/2022 a jul/2025, solar de set/2025 a ago/2026. Parte da diferença pode ser de período; o teste no Nordeste controla geografia, não tempo.}
\ressalva{A eólica se aglomera mais (mediana de 4 por nó contra 2 no solar), o que facilita detectar aglomeração --- mas a distância à carga tem faixa semelhante nas duas e não prevê o corte eólico.}
\ressalva{Regimes distintos nesta amostra: 8\% mediano em eólica e 26\% em solar; nas frotas, 15,9\% e 22,1\%. Origem é classificação do ONS, não inferência causal.}
}{\legXI}

%================================================================ FIG 12
\analisec{ANÁLISE 12 \quad · \quad Figs.\,10 e 11 + ONS \quad · \quad qualifica ambas}
{Robustez dos quatro resultados de rede que sobreviveram a Holm}
{Bootstrap, partição estrutural e reconstrução do grafo por outra fonte, aplicados aos quatro resultados. Um passa nos três critérios; o principal da Fig.\,10 passa em um.}
{../figures/fig12_robustez.pdf}
{
\bloco{Método}
O ONS publica a rede com \textbf{conectividade por número de barra}, impedância em
100\% das linhas e capacidade em MVA em 87\%, o que permite reconstruir o grafo por
via independente da camada geométrica usada nas Figs.\,10 e 11. Os três critérios são:
intervalo bootstrap que exclui zero, estabilidade entre as metades de uma partição
estrutural, e reprodução no grafo alternativo.

\bloco{Resultados}
\par\vspace{1pt}\scriptsize
\begin{tabular}{@{}lccc@{}}
\toprule
 & IC & estável & outro grafo \\
\midrule
eólica · grau do nó & sim & sim & sim \\
eólica · usinas no nó & sim & \textbf{não} & sim \\
solar · capacidade da usina & sim & sim & --- \\
solar · distância à carga & sim & \textbf{não} & \textbf{não} \\
\bottomrule
\end{tabular}
\par
\ressalva{--- não se aplica: a capacidade da usina não deriva do grafo. Dois critérios, e não três.}

Sobreviveram a Holm \textbf{quatro} resultados, e não três: a \textbf{capacidade da
usina} em solar --- $\rho = -0{,}38$, Holm $p = 0{,}023$, magnitude igual à da distância
à carga --- não constava do conjunto avaliado.

A distância à carga tem intervalo bootstrap de 0,12 a 0,61 --- \textbf{largura 0,49},
cerca do dobro dos demais. Partindo a amostra, o coeficiente cai de \textbf{0,62} para
\textbf{0,11} conforme a metade ($p = 0{,}025$), e a metade de maior coeficiente contém
\textbf{todas as 16 usinas do Sudeste}, o que reabre o confundimento regional que o
painel \textbf{d} da Fig.\,10 afastava.

A capacidade da usina passa nos dois critérios aplicáveis: IC de $-0{,}60$ a $-0{,}13$
e partição não rejeitada ($-0{,}18$ contra $-0{,}48$, $p = 0{,}21$). Dois critérios não
são três, e a diferença entre metades ainda reduz o coeficiente em dois terços:
\textbf{passar aqui é não ser rejeitado, não é ser confirmado}.

\bloco{Interpretação}
No grafo do ONS a distância dá $-0{,}14$ em km e $-0{,}06$ em reatância. Rodando o grafo
da EPE nas mesmas 28 usinas, ele também não apresenta associação: a diferença é de
amostra, o que não recupera o resultado --- identifica em que subconjunto ele ocorre.

A distância elétrica foi calculada, e a afirmação da Fig.\,10 de que ela não estaria
disponível é incorreta. Sua correlação com a distância em km é de \textbf{0,53}, de modo
que a substituição constitui teste independente. O resultado não muda.

\bloco{Revisão}
Os três critérios usam apenas dados já disponíveis na Fig.\,10 e poderiam ter sido
aplicados naquela análise; foram executados a partir da troca de grafo.

\lit{Proctor, Carleton e Sum (2023): variável derivada de sensoriamento --- como este
grafo --- propaga erro para o intervalo, e nenhuma imputação foi aplicada.}

\bloco{Limitações}
\ressalva{A partição não é aleatória, por construção: testa dependência a característica estrutural, não variabilidade amostral genérica.}
\ressalva{O grafo do ONS é exato onde cobre, mas exclui conexões abaixo de 230 kV --- para a amostra solar isso remove mais da metade e elimina um subsistema inteiro.}
\ressalva{Três critérios não esgotam robustez: faltam validação temporal, análise de influência ponto a ponto e sensibilidade à definição de centro de carga. Os limiares são convenção declarada, e o teste do IC foi corrigido de ``é positivo'' para ``exclui zero''.}
}{\legXII}

%================================================================ AUX
\analise{PÁGINA AUXILIAR \quad · \quad apoio de decisão, fora da série}
{Imagem comercial: resolução, cobertura e custo}
{Os 21 sensores de um fornecedor típico posicionados contra a curva de detectabilidade medida na Fig.\,2, e o custo de varredura na ausência de coordenada por instalação.}
{../figures/aux_geopera_capacidade.pdf}
{
\bloco{Resultados}
A transição de detectabilidade ocorre \textbf{entre 1\,m e 5\,m}. Abaixo de 1\,m,
99,5\% da MMGD é segmentável e os 18 sensores submétricos estão no platô. Acima, a
fração cai: 86,6\% a 2\,m, 7,5\% a 5,3\,m, 1,6\% no Sentinel-2. Os dois sensores
hiperespectrais do catálogo ficam acima do joelho da curva.

A faixa de 2\,m custa AUD 2--3/km$^2$ e entrega a mesma resolução do CBERS-4A WPM, que
é público. Entre 0,80\,m e 2\,m há apenas arquivo, sem programação de aquisição.

\bloco{Interpretação}
O painel \textbf{c} quantifica a limitação registrada na Fig.\,2. Com coordenada por
instalação, um chip de 0,01\,km$^2$ a 30\,cm custa AUD 0,18--0,28. Sem ela, a varredura
municipal de 1.000\,km$^2$ custa AUD 18 mil--28 mil --- cinco ordens de grandeza. O
custo é determinado pela ausência de localização no registro, não pela resolução
disponível.

A versão reduzida é executável: varrer 100--400\,km$^2$ da área ocupada de dois ou três
municípios rurais custa AUD 1.800--11.200 e permite auditar a \textbf{contagem
agregada} contra o registro municipal. A verificação por instalação na classe rural
permanece inviável.

As 361.622 instalações de \textbf{pessoa jurídica} têm CEP completo no registro,
geocodificável em nível de rua. Para esse segmento --- comércio e indústria, com
arranjos medianos de 66 e 92\,m$^2$ --- o chip dirigido é viável, com custo da ordem de
AUD 0,20 por instalação.

A análise energética por satélite em uso operacional emprega satélite
\textbf{meteorológico geoestacionário} --- índice de nuvem $\rightarrow$ GHI
$\rightarrow$ modelo de desempenho --- a 1--2\,km e 10 minutos. Ela mede a nuvem, não o
arranjo, e é a abordagem das Figs.\,5 e 6.

\bloco{Limitações}
\ressalva{Preços e sensores são os publicados pelo fornecedor, em AUD e sem tributo; mudam sem aviso.}
\ressalva{A profundidade de arquivo não é publicada por sensor, o que é decisivo porque todo desenho da série é longitudinal.}
\ressalva{Catálogo só óptico: sem SAR, que seria o complemento imune a nuvem.}
}{{\color{claro}Página auxiliar --- ver Fig.\,2 para a curva de detectabilidade.}}

%================================================================ APLICAÇÃO
\cab{APLICAÇÃO \quad · \quad exercício fora da série numerada}
{Produção e exposição ao corte de uma área nova em Jaíba}
{Encadeamento dos fatores medidos nas figuras anteriores para estimar quanto uma área nova junto a Sol do Cerrado produziria e quanto dessa energia o despacho absorveria.}

\begin{minipage}[t]{0.315\linewidth}\vspace{0pt}\small
\bloco{Método}
Todos os fatores são medidos em figuras anteriores; nenhum é suposto.
\par\vspace{2pt}\scriptsize
\begin{tabular}{@{}lll@{}}
\toprule
passo & fator & origem \\
\midrule
terreno $\rightarrow$ módulo & $\div$\,2,5 & Fig.\,3 \\
módulo $\rightarrow$ capacidade & $\div$\,5,12\,m²/kWp & Fig.\,2 \\
capacidade $\rightarrow$ potencial & $\times$\,0,229 de FC & Figs.\,5 e 6 \\
\bottomrule
\end{tabular}
\par\small\vspace{2pt}
O fator de capacidade potencial vem da razão de desempenho limpa de Sol do Cerrado
(0,931, calibrada em 279 dias sem restrição) e do GHI da NASA POWER sobre Jaíba
(5,908\,kWh/m²/dia).

\bloco{Validação da cadeia}
Aplicada à usina existente, a cadeia fecha: 681,3\,MW fiscalizados implicam 3,49\,km²
de módulo, e o caminho inverso devolve 681,3\,MW. \textbf{Isso não valida a cadeia}: é
divisão e multiplicação pela mesma constante de 5,12\,m²/kWp, e fecharia com qualquer
valor dela. A ida e volta confere a aritmética.

O elo não validado é o primeiro. A razão terreno/módulo de 2,5 é típica de projeto, não
medida nesta usina, e entra linearmente na capacidade: um erro de 20\% nela produz 20\%
de erro na capacidade. Medi-la exigiria delinear o arranjo por segmentação, o passo que
a Fig.\,7 deixa em aberto.

\bloco{Resultados}
Para 2025, o método prevê fator de capacidade de \textbf{0,229} e \textbf{1.368\,GWh}.
A usina entregou fator de \textbf{0,150} e \textbf{897\,GWh}.
\end{minipage}\hfill
\begin{minipage}[t]{0.315\linewidth}\vspace{0pt}\small
\bloco{Comparação com a apuração do ONS}
A diferença é de \textbf{34,4\%}. O ONS apurou, por outro caminho --- curva de
desempenho declarada da usina e apuração operativa de \textit{constrained-off} ---
uma taxa de \textbf{32,8\%} na mesma usina e ano. As duas diferem em \textbf{1,6
ponto}.

A independência entre as duas estimativas é \textbf{parcial}. A cadeia não usa a
geração de referência do operador, que é a parte externa. Usa a razão de desempenho
limpa de 0,931, calibrada nos dias que a \textit{própria apuração do ONS} marca como
sem restrição: o sinalizador de restrição é compartilhado, e o que se testa é a
estimativa de \textit{quanto}, não a de \textit{quando}.

Não é uma terceira estimativa: $1 - PR_{2025}/PR_{\text{limpo}}$ é o mesmo estimador da
Fig.\,6, que dá \textbf{33,5\%}; a diferença para os 34,4\% é de agregação do GHI.
\textbf{São duas estimativas, não três}, e o que esta página acrescenta à Fig.\,6 é o
passo área\,$\rightarrow$\,capacidade, que é o elo não validado.

\bloco{Cenários para uma área nova}
\par\vspace{1pt}\scriptsize
\begin{tabular}{@{}rrrrr@{}}
\toprule
AOI & capac. & potencial & a 31\% & a 62\% \\
km² & MW & GWh/ano & GWh & GWh \\
\midrule
 1 &    78 &   157 &   108 &    60 \\
 2 &   156 &   314 &   216 &   119 \\
 5 &   391 &   784 &   541 &   298 \\
10 &   781 & 1.569 & 1.082 &   596 \\
20 & 1.563 & 3.137 & 2.165 & 1.192 \\
\bottomrule
\end{tabular}
\par\small\vspace{2pt}
As duas últimas colunas separam taxa média de taxa marginal: \textbf{a taxa média não
se aplica a uma usina nova}. Novan e Wang (2024) mediram que a taxa \textbf{marginal}
--- quanto de uma usina nova é cortado --- é cerca de duas vezes a média, o que define
a coluna de 62\%.
\end{minipage}\hfill
\begin{minipage}[t]{0.315\linewidth}\vspace{0pt}\small
\bloco{Exposição ao corte no ponto}
O barramento de Jaíba tem \textbf{1.301\,MW} em três conjuntos --- Sol do Cerrado
(681), Jaíba V (500) e Jaíba 138\,kV (120) ---, todos cortados entre 24\% e 31\% na
janela de referência.

O corte no ponto é \textbf{89,5\% de origem sistêmica}: 61,8\% por razão energética,
16,2\% por confiabilidade e 11,5\% por indisponibilidade externa; 10,5\% tem origem
local. No agregado nacional a distribuição é 68,5\% energética, 22,0\% confiabilidade e
9,5\% indisponibilidade externa. A restrição não é de saturação do barramento, e sim de
oferta excedente no instante do corte.

O resultado é consistente com as Figs.\,8 a 10: térmica inflexível nas horas de corte,
corredor com folga e distância à demanda como preditor. Jaíba está a \textbf{605\,km de
rede} do centro de carga mais próximo, no terceiro quartil da amostra.

\bloco{Leitura não sustentada pelos dados}
A hipótese de que mais capacidade no mesmo nó produz mais corte não é sustentada: na
Fig.\,10 a capacidade renovável no nó dá $\rho = -0{,}24$ e não sobrevive a Holm, e
tampouco o grau, a tensão ou o número de usinas. O que sobrevive é a distância. A
leitura defensável é a de posição na rede, não a de saturação do ponto, e uma usina
nova ali herdaria essa posição.

\end{minipage}

\vspace{5pt}{\color{claro}\rule{\linewidth}{0.4pt}}\par\vspace{4pt}
\begin{minipage}[t]{0.485\linewidth}\vspace{0pt}\small
\bloco{Limitações}
\ressalva{O cálculo não usa limite de linha. O ONS publica capacidade operativa em MVA, com variantes sazonais, em 86,7\% das linhas; incorporá-la tornaria a estimativa de exposição ao corte mais precisa.}
\ressalva{A razão marginal de 2$\times$ é medida na Califórnia. Não há equivalente brasileiro publicado; é ordem de grandeza declarada, não fator calibrado nesta análise.}
\ressalva{A razão terreno/módulo de 2,5 é o elo não validado da cadeia; inclinação, rastreamento e relevo a deslocam.}
\ressalva{O fator de capacidade potencial usa o desempenho limpo de Sol do Cerrado; uma usina nova com outro módulo ou rastreamento teria outro.}
\end{minipage}\hfill
\begin{minipage}[t]{0.485\linewidth}\vspace{0pt}\small
\bloco{Escopo}
O resultado é uma estimativa \textbf{física e operativa}: quanto a área produziria e
quanto o despacho absorveria. Cobre dimensionamento, comparação de sítios e antecipação
de exposição ao corte. Uma decisão de investimento depende de contrato, outorga e
parecer de acesso, que não estão em dado aberto e não são tratados aqui.
\par\vspace{3pt}
{\color{claro}\scriptsize Dados-fonte em \texttt{data/processed/aplicacao\_*.csv},
gerados por \texttt{notebooks/\_aplicacao\_expansao.py}.}
\end{minipage}
\clearpage

%================================================================ FUNDAMENTAÇÃO
\cab{FUNDAMENTAÇÃO}{Posição na literatura}{O que já estava publicado, o que a série reproduz, o que ela acrescenta e o que mudou na regulação durante a análise.}

\begin{minipage}[t]{0.315\linewidth}\vspace{0pt}\footnotesize
\bloco{Comparação internacional}
\scriptsize
\begin{tabular}{@{}lll@{}}
\toprule
sistema & taxa & referência \\
\midrule
Brasil, solar (apurada) & \textbf{21,6\%} & \textit{esta série} \\
Brasil, eól.+solar 2025 & 20,7\% & ONS \\
Chile 2024 & 17\% & imprensa setorial \\
Brasil, eólica (apurada) & \textbf{15,9\%} & \textit{esta série} \\
Irlanda, eólica 2025 & 11,3\% & EirGrid \\
Espanha, jul/2025 & $\approx$11\% & recorde \\
Califórnia, solar & 4,3\% & Novan e Wang \\
\bottomrule
\end{tabular}
\par\vspace{3pt}\footnotesize
Os denominadores diferem, e a comparação com os sistemas estrangeiros é de ordem de
grandeza, não pareada. As \textbf{duas linhas brasileiras são pareadas}: ambas são razão
de energia --- cortada dividida por disponível --- na janela abr/2024--ago/2026, em que
solar e eólica têm apuração simultânea, sobre todos os conjuntos com apuração (87
solares, 178 eólicos). Numa única extração de hoje elas dão 22,1\% e 15,9\%; a linha da
solar preserva os 21,6\% da Fig.\,1, apurados em extração anterior.

As taxas brasileiras estão no topo da faixa internacional para as duas fontes, com a
solar acima do Chile e a eólica acima da Irlanda. A solar perde cerca de 1,4 vez o que a
eólica perde, diferença menor do que sugeriria comparar a taxa agregada de uma com a
mediana da outra. O que separa os dois regimes não é a magnitude do corte, e sim a
\textbf{composição da causa}: 91\% de origem sistêmica na solar contra 76\% na eólica,
medido na Fig.\,11.

Novan e Wang (2024) mostram que a taxa \textbf{marginal} --- quanto de uma usina
\textit{nova} será cortado --- é cerca de duas vezes a média; na Califórnia, média de
4,3\% contra marginal de 9\%. Caso a razão se mantenha, a exposição de um projeto
brasileiro entrando agora é superior aos 21,6\% medidos.

\bloco{Mecanismo e teste}
A competição entre mínimo técnico térmico e renovável variável é o problema central da
literatura chinesa de curtailment, e gerou resposta de política explícita: o programa de
\textit{flexibility retrofit}, com 16,35\,GW em demonstração em 2017.

O que não estava disponível era o teste com a decomposição de despacho publicada pelo
ONS, que separa geração por ordem de mérito de geração por inflexibilidade declarada,
razão elétrica, \textit{unit commitment} e segurança energética. As Figs.\,8 e 9 usam
esse campo, e a Fig.\,9 mostra que ele não sustenta atribuição local.
\end{minipage}\hfill
\begin{minipage}[t]{0.315\linewidth}\vspace{0pt}\footnotesize
\bloco{Resultados reproduzidos}
\textbf{MMGD como preditor do corte.} Vieira et al.\ (2026) identificaram a MMGD como
preditor mais forte do corte por demanda ($OR\approx9{,}8$ para solar), com a carga do
sistema atuando de forma protetiva. O gradiente por quintil de carga da Fig.\,8 reproduz
esse padrão de forma independente.

\textbf{Invisibilidade da geração atrás do medidor.} É problema documentado: a geração
distribuída não é observada pelo operador e eleva o erro de previsão de carga líquida em
30\% no nível domiciliar e até 250\% no agregado. A Fig.\,2 mede a outra face do mesmo
problema: ela também não é observável por satélite gratuito.

\textbf{Resolução necessária para telhado.} A literatura que detecta MMGD opera entre 0,1
e 0,3\,m. O inventário global de Kruitwagen et al.\ (2021), a referência com Sentinel-2,
cobre apenas instalações acima de 10\,kW.

\bloco{Acréscimos}
\textbf{Validação externa da referência do ONS.} A geração de referência é derivada de
curvas de desempenho da própria usina. Validá-la exige fonte independente --- feito no
lado eólico com anemometria de SCADA (Silva et al., 2026) e, na Fig.\,6, no lado solar
com irradiância de satélite, em 74 usinas.

\textbf{Resultado negativo sobre observabilidade.} A construção é detectável e datável
por Sentinel-2; a operação não é.

\textbf{Retirada de uma atribuição causal.} A Fig.\,9 mostra que a associação entre corte
e térmica inflexível não distingue causa local de fator comum, o que retira da Fig.\,8 a
atribuição causal ali formulada.

\textbf{Dependência do preditor à composição da causa.} A literatura documenta que as
duas causas existem (Bird et al.), mede a proporção entre elas (EirGrid) e sustenta que o
corte deve ser estimado por nó (Maji et al.). A Fig.\,11 testa se \textbf{a covariável que
prevê o corte muda conforme a composição da causa} --- preditor sistêmico onde o corte é
sistêmico, preditor local onde é local. Não foi localizado teste publicado equivalente; a
busca não foi sistemática.
\end{minipage}\hfill
\begin{minipage}[t]{0.315\linewidth}\vspace{0pt}\footnotesize
\bloco{Mudança regulatória durante a análise}
A questão que a Fig.\,3 deixou em aberto --- se o corte é ressarcido --- deixou de ser
hipotética.

A \textbf{Lei 15.269/2025} criou compensação por \textit{constrained-off} para eólicas e
solares, custeada por Encargos de Serviços do Sistema, com efeitos retroativos a
1\textsuperscript{o} de setembro de 2023 e condicionada à renúncia de ações judiciais e à
assinatura de termo de compromisso. A Portaria MME 140 regulamentou os cortes ocorridos
até 25 de novembro de 2025; os posteriores ficam com a ANEEL.

A primeira fase do processo reuniu manifestação de interesse de \textbf{1.539 usinas},
cerca de 53\,GW instalados. Estimativas setoriais apontam perda da ordem de
\textbf{R\$\,6,5 bilhões} em 2025 com \textit{constrained-off} de eólica e solar, mesma
ordem de grandeza da valoração a CMO da Fig.\,3 extrapolada para o parque.

Isso reforça a limitação daquela figura: CMO mede valor sistêmico, não ressarcimento. Com
a lei, os dois divergem por construção --- o gerador é compensado por energia cujo valor
sistêmico, no instante do corte, era próximo de zero.

\bloco{Qualidade das fontes}
As fontes não têm o mesmo peso. São revisadas por pares: Vieira et al.\ (2026), Novan e
Wang (2024), Kruitwagen et al.\ (2021), Yu et al.\ (2018), Jiang et al.\ (2021),
Alix-García e Millimet (2023). São \textit{working papers} ou pré-prints: Newbery (2024),
Proctor et al.\ (2023), Silva et al.\ (2026), Global Renewables Watch. São fontes
normativas: a Lei 15.269/2025 e a Portaria MME 140. São imprensa setorial as taxas de
Chile, Espanha e Europa e o valor de R\$\,6,5\,bi, usadas como ordem de grandeza e não
como evidência.

\end{minipage}
\clearpage
%================================================================ REFERÊNCIAS
\cab{REFERÊNCIAS}{Fontes citadas}{Bibliografia, bases de dados e atos normativos.}

\begin{minipage}[t]{0.485\linewidth}\vspace{0pt}\footnotesize
\bloco{Curtailment e despacho}
Vieira, G.T.T.; Silva, W.N.; Lourenço, L.F.N.; Monaro, R.M.; Salles, M.B.C.;
Almeida, C.F.M. \textit{Characterizing wind and solar curtailment in Brazil: an
evidence-based analysis of operational drivers.} Electric Power Systems Research,
2026. \par\vspace{2pt}
Novan, K.; Wang, Y. \textit{Estimates of the marginal curtailment rates for solar and
wind generation.} Journal of Environmental Economics and Management, 2024.
\par\vspace{2pt}
Newbery, D. \textit{Marginal curtailment of wind and solar PV: transmission
constraints, pricing and access regimes for efficient investment.} EPRG Working
Paper 2401, University of Cambridge, 2024. \par\vspace{2pt}
Silva, W.N. et al. \textit{An integrated methodology for assessing wind power
curtailment using anemometric measurements and operational data in the Brazilian
context.} Atmosphere, 2026. \par\vspace{2pt}
Bird, L.; Lew, D.; Milligan, M. et al. \textit{Wind and solar energy curtailment: a
review of international experience.} Renewable and Sustainable Energy Reviews, 65,
577--586, 2016. \par\vspace{2pt}
Maji, D.; Irwin, D.E.; Shenoy, P.J.; Sitaraman, R.K. \textit{A first look at
node-level curtailment of renewable energy and its implications.} ACM e-Energy '25,
Rotterdam, 2025. \par\vspace{2pt}
EirGrid. \textit{Annual renewable energy constraint and curtailment report}, edições
2013--2024. \par\vspace{2pt}
Clean Energy Ministerial. \textit{Thermal power plant flexibility.} Advanced Power
Plant Flexibility Campaign, 2018.

\bloco{Sensoriamento remoto de fotovoltaica}
Kruitwagen, L.; Story, K.T.; Friedrich, J.; Byers, L.; Skillman, S.; Hepburn, C.
\textit{A global inventory of photovoltaic solar energy generating units.} Nature,
598, 604--610, 2021. \par\vspace{2pt}
Yu, J.; Wang, Z.; Majumdar, A.; Rajagopal, R. \textit{DeepSolar: a machine learning
framework to efficiently construct a solar deployment database in the United
States.} Joule, 2(12), 2018. \par\vspace{2pt}
Jiang, H. et al. \textit{Multi-resolution dataset for photovoltaic panel segmentation
from satellite and aerial imagery.} Earth System Science Data, 13, 5389--5401, 2021.
\par\vspace{2pt}
Robinson, C. et al. \textit{Temporal cluster matching for change detection of
structures from satellite imagery.} Microsoft Research, 2021. \par\vspace{2pt}
Microsoft Research. \textit{Global Renewables Watch: a temporal dataset of solar and
wind energy derived from satellite imagery.} arXiv:2503.14860, 2025.
\end{minipage}\hfill
\begin{minipage}[t]{0.485\linewidth}\vspace{0pt}\footnotesize
\bloco{Método --- variáveis derivadas de sensoriamento}
Proctor, J.; Carleton, T.; Sum, S. \textit{Parameter recovery using remotely sensed
variables.} NBER Working Paper 30861, 2023. \par\vspace{2pt}
Alix-García, J.; Millimet, D.L. \textit{Remotely incorrect? Accounting for
nonclassical measurement error in satellite data on deforestation.} Journal of the
Association of Environmental and Resource Economists, 10(5), 1335--1367, 2023.
\par\vspace{2pt}
Callaway, B.; Sant'Anna, P.H.C. \textit{Difference-in-differences with multiple time
periods.} Journal of Econometrics, 225(2), 2021.

\bloco{Atos normativos}
Lei n\textsuperscript{o} 15.269, de 2025 --- compensação por \textit{constrained-off}
de fontes eólica e solar. \par\vspace{2pt}
Portaria MME n\textsuperscript{o} 140 --- procedimentos de compensação para cortes
entre 1/9/2023 e 25/11/2025. \par\vspace{2pt}
Lei n\textsuperscript{o} 14.300, de 2022 --- marco legal da microgeração e
minigeração distribuída.

\bloco{Bases de dados}
EPE, webmap institucional (serviço ArcGIS REST): camadas de linhas de transmissão e
subestações em operação, com geometria. Alternativa tabular: ANEEL, SIGET ---
\textit{Contrato Módulo Linha Transmissão Subestação Origem Subestação Destino}.
\par\vspace{2pt}
ONS, Portal de Dados Abertos (\texttt{dados.ons.org.br}): restrição por
\textit{constrained-off} fotovoltaico e eólico; geração térmica por motivo de
despacho; fator de capacidade; curva de carga; intercâmbio nacional; energia
armazenada por subsistema; CMO semi-horário; relação conjunto--usina.
\par\vspace{2pt}
ANEEL: Sistema de Informações de Geração (SIGA); relação de empreendimentos de
geração distribuída e informações técnicas fotovoltaicas. \par\vspace{2pt}
ESA/Copernicus: Sentinel-2 L2A, via catálogo STAC público da AWS
(\texttt{earth-search.aws.element84.com}). \par\vspace{2pt}
NASA POWER: irradiância diária (MERRA-2 e SYN1DEG). \par\vspace{2pt}
IBGE: malhas territoriais das unidades da federação.

\bloco{Fontes de imprensa setorial}
Taxas de \textit{curtailment} de Chile, Espanha e Europa e estimativa de perda
financeira no Brasil: \textit{pv magazine}, Rystad Energy, Canal Solar, EIA
\textit{Today in Energy}. Usadas como ordem de grandeza.
\end{minipage}
\clearpage


%================================================================ SÍNTESE
\cab{SÍNTESE}{Resultados estabelecidos, revisados e em aberto}{Balanço ao fim de doze análises.}

\begin{minipage}[t]{0.315\linewidth}\vspace{0pt}\small
\bloco{Estabelecido}
\textbf{O corte é de grande magnitude e origem sistêmica.} 21,6\% no parque com
apuração, dois terços por razão energética, incidindo sobre toda a janela solar.

\textbf{A referência do ONS é calibrada.} Viés mediano 1,02 em 74 usinas, sem
degradação onde o corte é mais intenso. O agregado da Fig.\,1 permanece válido.

\textbf{A perda de Sol do Cerrado é de rede, não de recurso.} Nove décimos da queda de
29\% no fator de capacidade não se explicam pela irradiância.

\textbf{A MMGD é sub-pixel.} 91\% abaixo de um pixel Sentinel-2; a verificação por
imagem exige 2\,m ou melhor e coordenada por instalação, que não existe no registro.

\textbf{A construção é detectável e datável; a operação não.} +34\% de brilho relativo
na construção ($p=1{,}2\times10^{-5}$) contra ausência de sinal na operação
($p=0{,}09$).

\textbf{O piso térmico inflexível está no Sudeste.} 1.739\,MWmed de inflexibilidade
pura, contra 36 no Nordeste.

\textbf{O grau do nó está associado ao corte eólico.} $\rho = +0{,}35$, único dos quatro
resultados de rede a passar em bootstrap, partição estrutural e troca de grafo
(Fig.\,12).

\textbf{O preditor acompanha a origem do corte.} Solar é 91\% sistêmico e responde à
distância à demanda; eólica tem 24\% de corte local e responde a aglomeração no nó
($\rho = +0{,}47$). A dissociação se mantém sob controle de geografia.

\textbf{O corte pode ser estimado sem a referência do operador.} Irradiância mais razão
de desempenho calibrada estimam 33,5\% de perda em Sol do Cerrado em 2025 (Fig.\,6) e
34,4\% com o passo de área acoplado, contra 32,8\% apurados pelo ONS --- 1,6 ponto de
diferença. A independência é \textbf{parcial}: dispensa a geração de referência, mas usa
a apuração do ONS para selecionar os dias de calibração.
\end{minipage}\hfill
\begin{minipage}[t]{0.315\linewidth}\vspace{0pt}\small
\bloco{Conclusões revisadas}
\textbf{``A referência do ONS superestima o corte''} (Fig.\,5). Artefato do
contrafactual: 2023 já era um ano com restrição. Revisto na Fig.\,6, que \textit{reforça}
a conclusão principal --- corte implícito de 33\%, não 26\%.

\textbf{``A térmica inflexível desloca a geração solar''} (Fig.\,8). A decomposição da
Fig.\,9 mostra que essa térmica se associa igualmente ao corte a dois mil quilômetros:
é marcador sistêmico, e o desenho não separa as hipóteses.

\textbf{``O Nordeste é cortado porque o corredor satura''.} Nas horas de corte alto o
fluxo NE$\rightarrow$SE é menor que nas demais.

\textbf{``A térmica não abre espaço para a solar''.} Não sobrevive ao controle de carga:
a carga líquida diária é plana.

\textbf{``Explicações locais do corte generalizam''} (Fig.\,3). Nenhuma das cinco
covariáveis testadas se reproduz no parque.

\textbf{``A distância à demanda é sobre posição na rede, não sobre tecnologia''}
(Fig.\,10). A Fig.\,11 restringe o enunciado a fonte de corte sistêmico, e a Fig.\,12
mede sua fragilidade: cai de 0,62 para 0,11 conforme a metade da amostra e não aparece
no grafo por barra.

\textbf{``Impedância e limite de linha não estão em dado aberto''} (Figs.\,3, 10, 11 e
aplicação). Estão, no pacote \texttt{linha-transmissao} do ONS: reatância em 100\% das
linhas e capacidade em MVA em 87\%. A afirmação incorreta aparecia em quatro pontos das
versões anteriores.

\textbf{``Detectar painéis por satélite audita política de MMGD''} --- a premissa
original, encerrada pela Fig.\,2.
\end{minipage}\hfill
\begin{minipage}[t]{0.315\linewidth}\vspace{0pt}\small
\bloco{Em aberto}
\textbf{Alocação do corte entre usinas, parcialmente respondida.} A Fig.\,10 identificou
duas covariáveis que sobrevivem à correção de multiplicidade: distância de rede à demanda
e capacidade da usina. Sob os critérios da Fig.\,12 é a \textbf{segunda} que resiste, e é
a que não tem explicação testada. A variância explicada é pequena, e a maior parte da
dispersão da taxa de corte permanece sem explicação.

\textbf{Causa do corte no Nordeste.} Nem piso térmico local, nem escoamento. Resta
ausência de demanda no destino, obtida por eliminação.

\textbf{Correspondência entre inflexibilidade declarada e mínimo técnico real.} É a
grandeza que uma auditoria exigiria, e não está publicada.

\bloco{Trabalho futuro, em ordem de retorno esperado}
\textbf{1. Reconstruir o corte anterior a abr/2024.} A maquinaria da Fig.\,6 já existe;
um modelo treinado nos intervalos sem restrição, aplicado retroativamente, viabiliza o
teste temporal da Fig.\,4. Há verdade de campo para validar.

\textbf{2. Medir o resíduo não explicado.} No nível de intervalo, a diferença de
desempenho entre um modelo só com observáveis sistêmicos e outro com identidade de usina
quantifica o que as Figs.\,10 e 11 deixaram sem explicação. A série eólica acrescenta
1,4\,GB e três anos de painel.

\textbf{3. Substituir NASA POWER por NSRDB.} De célula de $\sim$50\,km e passo diário
para 2\,km e 10 minutos, gratuito. Atende três limitações já registradas.

\textbf{4. Delinear o arranjo por segmentação.} Substitui os círculos de 400\,m da
Fig.\,7 e aumenta a precisão da datação.

\bloco{Regra de método para as Figs.\,1 e 2}
Qualquer saída de modelo que entre em inferência a jusante enviesa o ponto e o erro
padrão. A predição deve entrar como distribuição, não como valor pontual.
\end{minipage}

\vfill
{\color{claro}\rule{\linewidth}{0.4pt}}\par\vspace{3pt}
{\color{claro}\scriptsize
Todas as análises usam exclusivamente dados abertos e sem autenticação. Cada figura
exporta seus dados-fonte em CSV e passa por auditoria de tamanho mínimo de glifo e de
colisão de elementos antes de ser considerada pronta. Notebooks e scripts geradores
estão em \texttt{notebooks/}; figuras em \texttt{../figures/}; dados-fonte em
\texttt{data/processed/}.}

\end{document}
